
\documentclass{article}
\usepackage{colortbl}
\usepackage{makecell}
\usepackage{multirow}
\usepackage{supertabular}

\begin{document}

\newcounter{utterance}

\twocolumn

{ \footnotesize  \setcounter{utterance}{1}
\setlength{\tabcolsep}{0pt}
\begin{supertabular}{c@{$\;$}|p{.15\linewidth}@{}p{.15\linewidth}p{.15\linewidth}p{.15\linewidth}p{.15\linewidth}p{.15\linewidth}}

    \# & $\;$A & \multicolumn{4}{c}{Game Master} & $\;\:$B\\
    \hline 

    \theutterance \stepcounter{utterance}  

    & & \multicolumn{4}{p{0.6\linewidth}}{\cellcolor[rgb]{0.9,0.9,0.9}{%
	\makecell[{{p{\linewidth}}}]{% 
	  \tt {\tiny [A$\langle$GM]}  
	 State giocando a un gioco di negoziazione cooperativa in cui dovete accordarvi su come dividere una serie di oggetti tra voi e un altro giocatore.\\ \tt \\ \tt Regole:\\ \tt (a) A Lei e all'altro giocatore viene dato un insieme di oggetti. Ognuno di voi riceve anche una funzione di valore segreta, che rappresenta il valore di ciascun tipo di oggetto.\\ \tt (b) Si scambiano messaggi con l'altro giocatore per concordare chi si aggiudica gli oggetti. Potete inviare un massimo di 5 messaggi ciascuno, oppure terminare in anticipo facendo una proposta segreta in qualsiasi momento.\\ \tt (c) A ciascuno di voi viene chiesto di inviare una proposta segreta indicando gli oggetti che desiderate, formattata tra parentesi quadre come segue: "[Proposta: <numero> <nome oggetto>, <numero> <nome oggetto>, <...>]".\\ \tt (d) Se le vostre proposte sono complementari, cioè ci sono abbastanza oggetti per soddisfare entrambe le proposte, a ciascun giocatore viene assegnato un punteggio basato sulla somma dei valori degli oggetti ricevuti. In caso contrario, entrambi ricevono zero punti.\\ \tt (e) L'obiettivo è massimizzare la somma del proprio punteggio e di quello dell'altro giocatore.\\ \tt \\ \tt Cominciamo.\\ \tt \\ \tt L'insieme degli oggetti disponibili è:\\ \tt \\ \tt 1 carota, 4 borse, 1 cipolla\\ \tt \\ \tt La funzione valore segreta è:\\ \tt \\ \tt carota: 3, borsa: 1, cipolla: 3\\ \tt \\ \tt IMPORTANTE: i vostri messaggi, a meno che non si tratti di una proposta segreta, vengono trasmessi direttamente all'altro giocatore, quindi non includete alcuna risposta alle regole o testo di annuncio del vostro messaggio. Per fare una proposta segreta, utilizzate il formato indicato. Non utilizzare delle parentesi quadre quando si comunica all'altro giocatore, altrimenti verrà interpretata come una proposta segreta.\\ \tt  
	  } 
	   } 
	   } 
	 & \\ 
 

    \theutterance \stepcounter{utterance}  

    & \multicolumn{4}{p{0.6\linewidth}}{\cellcolor[rgb]{0.8,1,0.9}{%
	\makecell[{{p{\linewidth}}}]{% 
	  \tt {\tiny [A$\rangle$GM]}  
	 Ciao, vedo che ci sono diversi oggetti interessanti in gioco. La mia strategia iniziale sarebbe quella di capire cosa potrebbe essere più prezioso per te. Ci sono molte borse, quindi forse non sono il nostro obiettivo principale. La carota e la cipolla sembrano essere gli oggetti più rari e potenzialmente più preziosi. Che ne pensi? 
	  } 
	   } 
	   } 
	 & & \\ 
 

    \theutterance \stepcounter{utterance}  

    & & & \multicolumn{2}{p{0.3\linewidth}}{\cellcolor[rgb]{0.95,0.95,0.95}{%
	\makecell[{{p{\linewidth}}}]{% 
	  \tt {\tiny [GM$|$GM]}  
	 continue 
	  } 
	   } 
	   } 
	 & & \\ 
 

    \theutterance \stepcounter{utterance}  

    & & & \multicolumn{2}{p{0.3\linewidth}}{\cellcolor[rgb]{0.95,0.95,0.95}{%
	\makecell[{{p{\linewidth}}}]{% 
	  \tt {\tiny [GM$|$GM]}  
	 Ciao, vedo che ci sono diversi oggetti interessanti in gioco. La mia strategia iniziale sarebbe quella di capire cosa potrebbe essere più prezioso per te. Ci sono molte borse, quindi forse non sono il nostro obiettivo principale. La carota e la cipolla sembrano essere gli oggetti più rari e potenzialmente più preziosi. Che ne pensi? 
	  } 
	   } 
	   } 
	 & & \\ 
 

    \theutterance \stepcounter{utterance}  

    & & \multicolumn{4}{p{0.6\linewidth}}{\cellcolor[rgb]{0.9,0.9,0.9}{%
	\makecell[{{p{\linewidth}}}]{% 
	  \tt {\tiny [GM$\rangle$B]}  
	 State giocando a un gioco di negoziazione cooperativa in cui dovete accordarvi su come dividere una serie di oggetti tra voi e un altro giocatore.\\ \tt \\ \tt Regole:\\ \tt (a) A Lei e all'altro giocatore viene dato un insieme di oggetti. Ognuno di voi riceve anche una funzione di valore segreta, che rappresenta il valore di ciascun tipo di oggetto.\\ \tt (b) Si scambiano messaggi con l'altro giocatore per concordare chi si aggiudica gli oggetti. Potete inviare un massimo di 5 messaggi ciascuno, oppure terminare in anticipo facendo una proposta segreta in qualsiasi momento.\\ \tt (c) A ciascuno di voi viene chiesto di inviare una proposta segreta indicando gli oggetti che desiderate, formattata tra parentesi quadre come segue: "[Proposta: <numero> <nome oggetto>, <numero> <nome oggetto>, <...>]".\\ \tt (d) Se le vostre proposte sono complementari, cioè ci sono abbastanza oggetti per soddisfare entrambe le proposte, a ciascun giocatore viene assegnato un punteggio basato sulla somma dei valori degli oggetti ricevuti. In caso contrario, entrambi ricevono zero punti.\\ \tt (e) L'obiettivo è massimizzare la somma del proprio punteggio e di quello dell'altro giocatore.\\ \tt \\ \tt Cominciamo.\\ \tt \\ \tt L'insieme degli oggetti disponibili è:\\ \tt \\ \tt 1 carota, 4 borse, 1 cipolla\\ \tt \\ \tt La funzione valore segreta è:\\ \tt \\ \tt carota: 0, borsa: 1, cipolla: 6\\ \tt \\ \tt IMPORTANTE: i vostri messaggi, a meno che non si tratti di una proposta segreta, vengono trasmessi direttamente all'altro giocatore, quindi non includete alcuna risposta alle regole o testo di annuncio del vostro messaggio. Per fare una proposta segreta, utilizzate il formato indicato. Non utilizzare delle parentesi quadre quando si comunica all'altro giocatore, altrimenti verrà interpretata come una proposta segreta.\\ \tt \\ \tt \\ \tt Ciao, vedo che ci sono diversi oggetti interessanti in gioco. La mia strategia iniziale sarebbe quella di capire cosa potrebbe essere più prezioso per te. Ci sono molte borse, quindi forse non sono il nostro obiettivo principale. La carota e la cipolla sembrano essere gli oggetti più rari e potenzialmente più preziosi. Che ne pensi? 
	  } 
	   } 
	   } 
	 & \\ 
 

    \theutterance \stepcounter{utterance}  

    & & & \multicolumn{4}{p{0.6\linewidth}}{\cellcolor[rgb]{1,0.85,0.72}{%
	\makecell[{{p{\linewidth}}}]{% 
	  \tt {\tiny [GM$\langle$B]}  
	 Sì, concordo che la carota e la cipolla potrebbero essere gli oggetti più preziosi, data la loro rarità. Tuttavia, non so ancora quale dei due potrebbe essere più importante per te. Penso che le borse potrebbero essere utili per entrambi, ma forse potremmo dividerle in modo da garantire a ciascuno un certo numero di esse. La vera domanda è: quanto valore attribuisci alla cipolla rispetto alla carota? 
	  } 
	   } 
	   } 
	  \\ 
 

    \theutterance \stepcounter{utterance}  

    & & & \multicolumn{2}{p{0.3\linewidth}}{\cellcolor[rgb]{0.95,0.95,0.95}{%
	\makecell[{{p{\linewidth}}}]{% 
	  \tt {\tiny [GM$|$GM]}  
	 continue 
	  } 
	   } 
	   } 
	 & & \\ 
 

    \theutterance \stepcounter{utterance}  

    & & & \multicolumn{2}{p{0.3\linewidth}}{\cellcolor[rgb]{0.95,0.95,0.95}{%
	\makecell[{{p{\linewidth}}}]{% 
	  \tt {\tiny [GM$|$GM]}  
	 Sì, concordo che la carota e la cipolla potrebbero essere gli oggetti più preziosi, data la loro rarità. Tuttavia, non so ancora quale dei due potrebbe essere più importante per te. Penso che le borse potrebbero essere utili per entrambi, ma forse potremmo dividerle in modo da garantire a ciascuno un certo numero di esse. La vera domanda è: quanto valore attribuisci alla cipolla rispetto alla carota? 
	  } 
	   } 
	   } 
	 & & \\ 
 

    \theutterance \stepcounter{utterance}  

    & & \multicolumn{4}{p{0.6\linewidth}}{\cellcolor[rgb]{0.9,0.9,0.9}{%
	\makecell[{{p{\linewidth}}}]{% 
	  \tt {\tiny [A$\langle$GM]}  
	 Sì, concordo che la carota e la cipolla potrebbero essere gli oggetti più preziosi, data la loro rarità. Tuttavia, non so ancora quale dei due potrebbe essere più importante per te. Penso che le borse potrebbero essere utili per entrambi, ma forse potremmo dividerle in modo da garantire a ciascuno un certo numero di esse. La vera domanda è: quanto valore attribuisci alla cipolla rispetto alla carota? 
	  } 
	   } 
	   } 
	 & \\ 
 

    \theutterance \stepcounter{utterance}  

    & \multicolumn{4}{p{0.6\linewidth}}{\cellcolor[rgb]{0.8,1,0.9}{%
	\makecell[{{p{\linewidth}}}]{% 
	  \tt {\tiny [A$\rangle$GM]}  
	 La tua osservazione sulle borse è interessante, potrebbero essere utili per entrambi. Per quanto riguarda la cipolla e la carota, posso dire che hanno un certo valore per me, ma non sono disposto a rivelare ancora il mio ordine di preferenza. Tuttavia, potrei essere disposto a fare una proposta che tenga conto dei tuoi interessi. Come ti sentiresti se ci dividessimo le borse in modo equo e poi discutessimo sulla carota e la cipolla? Saresti disposto a lasciarmi la carota se ti dessi la maggior parte delle borse? 
	  } 
	   } 
	   } 
	 & & \\ 
 

    \theutterance \stepcounter{utterance}  

    & & & \multicolumn{2}{p{0.3\linewidth}}{\cellcolor[rgb]{0.95,0.95,0.95}{%
	\makecell[{{p{\linewidth}}}]{% 
	  \tt {\tiny [GM$|$GM]}  
	 continue 
	  } 
	   } 
	   } 
	 & & \\ 
 

    \theutterance \stepcounter{utterance}  

    & & & \multicolumn{2}{p{0.3\linewidth}}{\cellcolor[rgb]{0.95,0.95,0.95}{%
	\makecell[{{p{\linewidth}}}]{% 
	  \tt {\tiny [GM$|$GM]}  
	 La tua osservazione sulle borse è interessante, potrebbero essere utili per entrambi. Per quanto riguarda la cipolla e la carota, posso dire che hanno un certo valore per me, ma non sono disposto a rivelare ancora il mio ordine di preferenza. Tuttavia, potrei essere disposto a fare una proposta che tenga conto dei tuoi interessi. Come ti sentiresti se ci dividessimo le borse in modo equo e poi discutessimo sulla carota e la cipolla? Saresti disposto a lasciarmi la carota se ti dessi la maggior parte delle borse? 
	  } 
	   } 
	   } 
	 & & \\ 
 

    \theutterance \stepcounter{utterance}  

    & & \multicolumn{4}{p{0.6\linewidth}}{\cellcolor[rgb]{0.9,0.9,0.9}{%
	\makecell[{{p{\linewidth}}}]{% 
	  \tt {\tiny [GM$\rangle$B]}  
	 La tua osservazione sulle borse è interessante, potrebbero essere utili per entrambi. Per quanto riguarda la cipolla e la carota, posso dire che hanno un certo valore per me, ma non sono disposto a rivelare ancora il mio ordine di preferenza. Tuttavia, potrei essere disposto a fare una proposta che tenga conto dei tuoi interessi. Come ti sentiresti se ci dividessimo le borse in modo equo e poi discutessimo sulla carota e la cipolla? Saresti disposto a lasciarmi la carota se ti dessi la maggior parte delle borse? 
	  } 
	   } 
	   } 
	 & \\ 
 

    \theutterance \stepcounter{utterance}  

    & & & \multicolumn{4}{p{0.6\linewidth}}{\cellcolor[rgb]{1,0.85,0.72}{%
	\makecell[{{p{\linewidth}}}]{% 
	  \tt {\tiny [GM$\langle$B]}  
	 Mi sembra una proposta interessante. Dividere le borse in modo equo potrebbe essere un buon punto di partenza. Per quanto riguarda la carota e la cipolla, sono disposto a considerare la tua proposta, ma devo ammettere che la cipolla mi sembra particolarmente interessante. Lasciarti la carota in cambio della maggior parte delle borse potrebbe essere una possibilità, ma dipenderebbe da quante borse stiamo parlando. Se potessi avere almeno 3 borse, potrei essere disposto a lasciarti la carota e forse anche discutere sulla cipolla. Che ne pensi? 
	  } 
	   } 
	   } 
	  \\ 
 

    \theutterance \stepcounter{utterance}  

    & & & \multicolumn{2}{p{0.3\linewidth}}{\cellcolor[rgb]{0.95,0.95,0.95}{%
	\makecell[{{p{\linewidth}}}]{% 
	  \tt {\tiny [GM$|$GM]}  
	 continue 
	  } 
	   } 
	   } 
	 & & \\ 
 

    \theutterance \stepcounter{utterance}  

    & & & \multicolumn{2}{p{0.3\linewidth}}{\cellcolor[rgb]{0.95,0.95,0.95}{%
	\makecell[{{p{\linewidth}}}]{% 
	  \tt {\tiny [GM$|$GM]}  
	 Mi sembra una proposta interessante. Dividere le borse in modo equo potrebbe essere un buon punto di partenza. Per quanto riguarda la carota e la cipolla, sono disposto a considerare la tua proposta, ma devo ammettere che la cipolla mi sembra particolarmente interessante. Lasciarti la carota in cambio della maggior parte delle borse potrebbe essere una possibilità, ma dipenderebbe da quante borse stiamo parlando. Se potessi avere almeno 3 borse, potrei essere disposto a lasciarti la carota e forse anche discutere sulla cipolla. Che ne pensi? 
	  } 
	   } 
	   } 
	 & & \\ 
 

    \theutterance \stepcounter{utterance}  

    & & \multicolumn{4}{p{0.6\linewidth}}{\cellcolor[rgb]{0.9,0.9,0.9}{%
	\makecell[{{p{\linewidth}}}]{% 
	  \tt {\tiny [A$\langle$GM]}  
	 Mi sembra una proposta interessante. Dividere le borse in modo equo potrebbe essere un buon punto di partenza. Per quanto riguarda la carota e la cipolla, sono disposto a considerare la tua proposta, ma devo ammettere che la cipolla mi sembra particolarmente interessante. Lasciarti la carota in cambio della maggior parte delle borse potrebbe essere una possibilità, ma dipenderebbe da quante borse stiamo parlando. Se potessi avere almeno 3 borse, potrei essere disposto a lasciarti la carota e forse anche discutere sulla cipolla. Che ne pensi? 
	  } 
	   } 
	   } 
	 & \\ 
 

    \theutterance \stepcounter{utterance}  

    & \multicolumn{4}{p{0.6\linewidth}}{\cellcolor[rgb]{0.8,1,0.9}{%
	\makecell[{{p{\linewidth}}}]{% 
	  \tt {\tiny [A$\rangle$GM]}  
	 Sembra che stiamo facendo progressi. Io sono disposto a darti 3 borse se lasci la carota a me. Tuttavia, per quanto riguarda la cipolla, sono un po' esitante a rinunciarci. Forse potremmo trovare un compromesso anche lì, ma vorrei sapere cosa pensi di un'idea: se ti do 3 borse e lascio che tu abbia la cipolla, tu lasci la carota a me e magari mi dai una delle borse in più che ti ho dato, così da trovare un equilibrio? [Proposta: 1 carota, 1 borsa] 
	  } 
	   } 
	   } 
	 & & \\ 
 

    \theutterance \stepcounter{utterance}  

    & & & \multicolumn{2}{p{0.3\linewidth}}{\cellcolor[rgb]{0.95,0.95,0.95}{%
	\makecell[{{p{\linewidth}}}]{% 
	  \tt {\tiny [GM$|$GM]}  
	 proposal 
	  } 
	   } 
	   } 
	 & & \\ 
 

    \theutterance \stepcounter{utterance}  

    & & & \multicolumn{2}{p{0.3\linewidth}}{\cellcolor[rgb]{0.95,0.95,0.95}{%
	\makecell[{{p{\linewidth}}}]{% 
	  \tt {\tiny [GM$|$GM]}  
	 [1, 1, 0] 
	  } 
	   } 
	   } 
	 & & \\ 
 

    \theutterance \stepcounter{utterance}  

    & & \multicolumn{4}{p{0.6\linewidth}}{\cellcolor[rgb]{0.9,0.9,0.9}{%
	\makecell[{{p{\linewidth}}}]{% 
	  \tt {\tiny [GM$\rangle$B]}  
	 L'altro giocatore ha presentato una proposta segreta. Ora tocca a Lei presentare una sua proposta.\\ \tt  
	  } 
	   } 
	   } 
	 & \\ 
 

    \theutterance \stepcounter{utterance}  

    & & & \multicolumn{4}{p{0.6\linewidth}}{\cellcolor[rgb]{1,0.85,0.72}{%
	\makecell[{{p{\linewidth}}}]{% 
	  \tt {\tiny [GM$\langle$B]}  
	 [Proposta: 2 borse, 1 carota] 
	  } 
	   } 
	   } 
	  \\ 
 

    \theutterance \stepcounter{utterance}  

    & & & \multicolumn{2}{p{0.3\linewidth}}{\cellcolor[rgb]{0.95,0.95,0.95}{%
	\makecell[{{p{\linewidth}}}]{% 
	  \tt {\tiny [GM$|$GM]}  
	 proposal 
	  } 
	   } 
	   } 
	 & & \\ 
 

    \theutterance \stepcounter{utterance}  

    & & & \multicolumn{2}{p{0.3\linewidth}}{\cellcolor[rgb]{0.95,0.95,0.95}{%
	\makecell[{{p{\linewidth}}}]{% 
	  \tt {\tiny [GM$|$GM]}  
	 [1, 2, 0] 
	  } 
	   } 
	   } 
	 & & \\ 
 

    \theutterance \stepcounter{utterance}  

    & & & \multicolumn{2}{p{0.3\linewidth}}{\cellcolor[rgb]{0.95,0.95,0.95}{%
	\makecell[{{p{\linewidth}}}]{% 
	  \tt {\tiny [GM$|$GM]}  
	 [[1, 1, 0], [1, 2, 0]] 
	  } 
	   } 
	   } 
	 & & \\ 
 

\end{supertabular}
}

\end{document}
